\documentclass[a4paper,12pt]{article}

\usepackage[T1]{fontenc}
\usepackage[italian]{babel}
\usepackage{graphicx}
\usepackage{geometry}
\usepackage{tabularx}
\usepackage[table]{xcolor}
\usepackage[most]{tcolorbox}
\usepackage{fancyhdr}
\usepackage[hidelinks, pdfnewwindow]{hyperref}

\newcommand{\groupmail}{alittlebyte.swe@gmail.com}
\newcommand{\grouplogo}{../../../.github/site-src/assets/logo_alittlebyte.png}
\newcommand{\groupsymbol}{../../../.github/site-src/assets/solo_logo.png}

\geometry{
    left=2.5cm,
    right=2.5cm,
    top=2.5cm,
    bottom=2.5cm,
    headheight=331pt,
    includeheadfoot
}

\pagestyle{fancy}
\fancyhf{}

\renewcommand{\headrulewidth}{0pt}
\fancyhead{}

\renewcommand{\footrulewidth}{0.4pt}
\fancyfoot[L]{\includegraphics[height=0.6cm]{\groupsymbol}\hspace{0.45em}aLittleByte}
    \fancyfoot[C]{\thepage}
\fancyfoot[R]{Verbale 2026-05-26}

\fancypagestyle{firstpage}{
    \fancyhf{}
    \renewcommand{\headrulewidth}{0pt}
    \renewcommand{\footrulewidth}{0.4pt}
    \fancyhead[C]{%
        \makebox[\textwidth][c]{%
            \begin{tabular}{c}
                \includegraphics[height=10cm]{\grouplogo}\\[1ex]
                \small\texttt{\groupmail}
            \end{tabular}%
        }
    }
    \fancyfoot[L]{\includegraphics[height=0.6cm]{\groupsymbol}\hspace{0.45em}aLittleByte}
        \fancyfoot[C]{\thepage}
    \fancyfoot[R]{Verbale 2026-05-26}
}

\providecommand{\glslink}[2]{\href{https://alittlebyte-19.github.io/Documentazione/glossario.html\#gls-#1}{\underline{#2}\textsuperscript{\scriptsize G}}}

\begin{document}

\thispagestyle{firstpage}

\vspace*{0.5cm}

\begin{center}
    {\LARGE \textbf{Verbale Interno - 2026-05-26}}
\end{center}

\vspace{1.5cm}

\begin{center}
    \renewcommand{\arraystretch}{1.3}
    \begin{tcolorbox}[
        width=0.92\textwidth,
        colback=gray!6,
        colframe=gray!45,
        boxrule=0.5pt,
        arc=2mm,
        boxsep=0pt,
        left=8pt, right=8pt, top=8pt, bottom=8pt
    ]
        \begin{tabularx}{\linewidth}{>{\bfseries}p{3.5cm} X}
            Gruppo      & aLittleByte (gruppo 19) \\
            Redattore   & Diego Belluz \\
            Verificatori & Angela Maule\\
            Destinatari & Prof. Tullio Vardanega, Prof. Riccardo Cardin \\
            Data        & 2026-05-26\\
        \end{tabularx}
    \end{tcolorbox}
\end{center}

\newpage

\newgeometry{left=2.5cm,right=2.5cm,top=2.5cm,bottom=2.5cm,includefoot}
\tableofcontents
\newpage
\section{Informazioni Generali}
\section*{Specifiche Documento}
\hrule
\vspace{1cm}

\begin{tabular}{>{\bfseries}p{5cm} | p{10cm}}
Luogo Riunione & \glslink{discord}{Discord} \\[6pt]

Orario (Inizio - Fine) & 10:00 -- 11:30 \\[6pt]

\glslink{redattore}{Redattore} &
  D. Belluz\\[6pt]

\glslink{amministratore}{Amministratore} & 
L. Soligo\\[6pt]

\glslink{verificatore}{Verificatore} &
A. Maule\\[6pt]

Partecipanti &
  A. Oloni\\ &
  L. Soligo\\&
  A. Gastaldello\\&
  E. Bellet\\&
  D. Belluz\\&
  A. Maule\\ [6pt]

Assentine motivazione & 
V. Y. Kaneda Fernandes - Impegni lavorativi\\


\end{tabular}


\section{Ordine del Giorno}
    \begin{itemize}
    \item Analisi del \glslink{preventivo}{preventivo} a finire relativo al quarto \glslink{sprint}{sprint} e valutazione delle varianze.
    \item \glslink{pianificazione}{Pianificazione} della consegna e definizione della data per \\ l'\glslink{rtb-requirements-and-technology-baseline}{RTB (Requirements and Technology Baseline)}.
    \item Verifica dello stato di avanzamento della documentazione (\glslink{norme-di-progetto}{Norme di Progetto}, \glslink{piano-di-progetto}{Piano di Progetto}, \glslink{analisi-dei-requisiti-adr}{Analisi dei Requisiti} e \glslink{piano-di-qualifica}{Piano di Qualifica}).
    \item Pubblicazione e aggiornamento dei documenti sul sito web del gruppo.
    \item Rotazione dei ruoli operativi e assegnazione dei nuovi incarichi.
\end{itemize}


\section{Attività svolte}
\subsection{Analisi del preventivo a finire relativo al quarto sprint e valutazione delle varianze}
Il team ha esaminato approfonditamente il \glslink{preventivo}{preventivo} a finire relativo al quarto \glslink{sprint}{sprint}, rilevando un trend di varianza crescente rispetto alle stime preventive. 
È stata presa piena consapevolezza dell'attuale scostamento, stabilendo la necessità di intervenire tempestivamente con un ribilanciamento strutturato del monte ore 
lavorative e dei costi associati per rientrare nei parametri preventivati.

\subsection{Pianificazione della consegna e definizione della data per l'RTB e Verifica dello stato di avanzamento della documentazione}
È stata concordata e formalizzata la data di consegna per la revisione \glslink{rtb-requirements-and-technology-baseline}{RTB}, fissata per il 29/05/2026. Durante il controllo ispettivo sullo stato dei documenti, 
il team ha confermato il completamento definitivo del "\glslink{piano-di-progetto}{Piano di Progetto}" e dell'"\glslink{analisi-dei-requisiti-adr}{Analisi dei Requisiti}". È emersa invece l'esigenza di effettuare una revisione mirata sulle 
"\glslink{norme-di-progetto}{Norme di Progetto}", al fine di redigere e integrare le sezioni attualmente mancanti, con priorità assoluta per il capitolo relativo alle metriche di qualità. È stato inoltre 
confermato l'impegno a ultimare la stesura del "\glslink{piano-di-qualifica}{Piano di Qualifica}" in vista della scadenza prefissata.


\subsection{Pubblicazione e aggiornamento dei documenti sul sito web del gruppo.}
È stata ribadita la necessità di completare il caricamento di tutta la documentazione aggiornata e validata all'interno del sito web del gruppo.

\subsection{Rotazione dei ruoli operativi e assegnazione dei nuovi incarichi.}
Infine, si è proceduto alla consueta rotazione dei ruoli di progetto: Diego Belluz assume formalmente la carica di \glslink{responsabile}{Responsabile} e viene incaricato della redazione del presente verbale, 
nonché della compilazione del \glslink{diario-di-bordo}{Diario di Bordo}, la cui pubblicazione è calendarizzata per il 29/05/2026.

\section{To-Do List}
\begin{table}[h]
    \centering
    \begin{tabular}{|p{4.5cm}|p{5cm}|l|}
    \hline
        \rowcolor{lightgray}\textit{\textbf{Persona Incaricata}}  & \textbf{Task Assegnata} &\textbf{Link \glslink{issue}{Issue}} \\
    \hline
        \textit{A. Oloni} & Aggiornamento \glslink{glossario}{Glossario} & \href{https://github.com/aLittleByte-19/Documentazione/issues/53}{\#53}\\
    \hline 
        \textit{V. Y. K. Fernandes} & Aggiornamento sito web & \href{https://github.com/aLittleByte-19/Documentazione/issues/103}{\#103}\\
    \hline
        \textit{Leonardo Soligo} &  Continuazione della stesura di \glslink{norme-di-progetto}{Norme di Progetto} & \href{https://github.com/aLittleByte-19/Documentazione/issues/52}{\#52}\\
    \hline
        \textit{Diego Belluz} & \glslink{diario-di-bordo}{Diario di Bordo} 29/05/2026 & \href{https://github.com/aLittleByte-19/Documentazione/issues/105}{\#105}\\
    \hline
        \textit{Diego Belluz} & Continuzione stesura \glslink{piano-di-qualifica}{piano di qualifica} & \href{https://github.com/aLittleByte-19/Documentazione/issues/92}{\#92}\\
    \hline
        \textit{Diego Belluz} & Stesura \glslink{verbale-interno}{Verbale Interno} 26/05/2026 & \href{https://github.com/aLittleByte-19/Documentazione/issues/102}{\#102}\\
    \hline
    \end{tabular}
    \label{tab:placeholder}
\end{table}

\end{document}